% Existing FormalizingGMT declarations at the pinned revision.
% The declarations and proof bodies are formalized upstream. No claim is made
% that the full downstream project has been rebuilt on this workstation.
\section{Existing FormalizingGMT results in the graph}
Green statement/proof markers in this section mean that the corresponding
definitions and proof bodies already exist in the pinned FormalizingGMT source.
They do not mean that the new Preiss adapters are proved, or that this workstation
has completed a fresh build of every imported module. The edges expose the
mathematical dependencies relevant here, not an exhaustive kernel proof-term graph.
Each node's Lean link opens a declaration/source viewer with its exact source location.

\begin{definition}[GMT: dimensional densities]
\label{def:density-api}
\lean{dimensional_density_ratio,dimensional_lower_density,dimensional_upper_density,HasDensity}
\leanok
For an outer measure, the density ratio is $\mu(\overline B(x,r))/(2r)^s$; lower and upper density use $r\downarrow0$. Density existence is convergence in $[0,\infty]$.
\end{definition}

\begin{lemma}[GMT: lower versus upper density]
\label{gmt:density-order}
\lean{lower_le_upper_density}
\leanok
\uses{def:density-api}
For every exponent and point, $\Theta_*^s(\mu,x)\leq\Theta^{*s}(\mu,x)$.
\end{lemma}
\begin{proof}[Formalized in FormalizingGMT]
\leanok
Import the existing declaration.
\end{proof}

\begin{lemma}[GMT: equal densities give a limit]
\label{gmt:density-limit}
\lean{density_exists_of_lower_eq_upper}
\leanok
\uses{def:density-api}
Equality of lower and upper density implies \texttt{HasDensity}. Positivity and finiteness remain separate hypotheses.
\end{lemma}
\begin{proof}[Formalized in FormalizingGMT]
\leanok
Import the existing declaration.
\end{proof}

\begin{lemma}[GMT: eventual small-ball bounds]
\label{gmt:density-bounds}
\lean{eventually_gt_of_lower_density_gt,eventually_lt_of_upper_density_lt,dimensional_density_ratio_closedBall}
\leanok
\uses{def:density-api}
Strict bounds below the lower density or above the upper density become eventual bounds for all sufficiently small positive radii. For nonnegative radii, the extended-ball expression agrees with the ordinary closed-ball ratio.
\end{lemma}
\begin{proof}[Formalized in FormalizingGMT]
\leanok
Import the existing declaration.
\end{proof}

\begin{theorem}[GMT: fine Vitali covering variant]
\label{gmt:vitali}
\lean{vitali_variant_classical}
\leanok
In a second-countable metric space, a fine centred family of positive-radius balls with bounded radii has a countable disjoint subfamily. After removing any finite subfamily of the original family, the remaining set is covered by the fivefold enlargements of the selected balls not removed.
\end{theorem}
\begin{proof}[Formalized in FormalizingGMT]
\leanok
Import the existing declaration.
\end{proof}

\begin{theorem}[GMT: finite restrictions are regular]
\label{gmt:finite-restriction}
\lean{BorelRegularOuterMeasure.restrict_isRadon}
\leanok
On a locally compact, Hausdorff, second-countable Borel space, $\mu E<\infty$ implies that $\mu|E$ is regular. The source does not require $E$ to be Borel.
\end{theorem}
\begin{proof}[Formalized in FormalizingGMT]
\leanok
Import the existing declaration.
\end{proof}

\begin{definition}[GMT: Hausdorff content]
\label{gmt:hausdorff-content}
\lean{hausdorffContent,hausdorffContentInfty}
\leanok
Hausdorff content is the infimum of countable covering diameter sums, with or without a diameter cutoff. For positive exponents, empty cover elements have zero contribution; the cutoff version uses an explicit nonempty guard.
\end{definition}

\begin{lemma}[GMT: content comparison and stabilization]
\label{gmt:content-comparison}
\lean{hausdorffContentInfty_le_hausdorffContent,hausdorffContent_le_hausdorffMeasure,hausdorffContent_eq_hausdorffContentInfty_of_ediam_le}
\leanok
\uses{gmt:hausdorff-content}
Unrestricted content is bounded above by cutoff content, which is bounded above by Hausdorff measure. For positive exponent and a positive cutoff at least the diameter of a set, cutoff content equals unrestricted content.
\end{lemma}
\begin{proof}[Formalized in FormalizingGMT]
\leanok
Import the existing declaration.
\end{proof}

\begin{theorem}[GMT: finite Hausdorff restrictions]
\label{gmt:hausdorff-regular}
\lean{HausdorffRestrict.toRadonOuterMeasure}
\leanok
\uses{gmt:hausdorff-content}
On a locally compact, second-countable Hausdorff metric space, restriction of $\mathcal H^s$ to a set of finite $\mathcal H^s$ measure is regular; use the source's Caratheodory-measurability and $s\geq0$ interface.
\end{theorem}
\begin{proof}[Formalized in FormalizingGMT]
\leanok
\uses{gmt:finite-restriction}
Import the existing declaration.
\end{proof}

\begin{theorem}[GMT: positive upper content density inside]
\label{gmt:content-density}
\lean{hausdorffContentInfty_upperDensity_ge_ae_mem}
\leanok
\uses{def:density-api,gmt:hausdorff-content}
For $s\geq0$ and $\mathcal H^s(E)<\infty$ in a metric space, the upper diameter-normalized density of restricted Hausdorff content is at least $2^{-s}$ at $\mathcal H^s$-almost every point of $E$.
\end{theorem}
\begin{proof}[Formalized in FormalizingGMT]
\leanok
\uses{gmt:content-comparison}
Import the existing declaration.
\end{proof}

\begin{theorem}[GMT: positive upper Hausdorff density inside]
\label{gmt:hausdorff-density-lower}
\lean{hausdorffMeasure_upperDensity_ge_ae_mem}
\leanok
\uses{def:density-api}
For $s\geq0$ and $\mathcal H^s(E)<\infty$, $\Theta^{*s}(\mathcal H^s|E,x)\geq2^{-s}$ at $\mathcal H^s$-almost every $x\in E$. This is a lower bound on the UPPER density, not on the lower density.
\end{theorem}
\begin{proof}[Formalized in FormalizingGMT]
\leanok
\uses{gmt:content-density,gmt:content-comparison}
Import the existing declaration.
\end{proof}

\begin{theorem}[GMT: upper Hausdorff density at most one]
\label{gmt:hausdorff-density-upper}
\lean{hausdorffMeasure_upperDensity_le_one_ae_mem}
\leanok
\uses{def:density-api}
If $E$ is Caratheodory measurable with finite $\mathcal H^s$ measure, $s\geq0$, and the metric space is locally compact and second countable, then $\Theta^{*s}(\mathcal H^s|E,x)\leq1$ for $\mathcal H^s$-almost every $x\in E$.
\end{theorem}
\begin{proof}[Formalized in FormalizingGMT]
\leanok
\uses{gmt:hausdorff-regular}
Import the existing declaration.
\end{proof}

\begin{theorem}[GMT: zero upper density outside]
\label{gmt:hausdorff-density-outside}
\lean{hausdorffMeasure_upperDensity_eq_zero_ae_notMem}
\leanok
\uses{def:density-api}
Under the same local compactness, second countability, measurability and finite-measure hypotheses, the upper $s$-density of $\mathcal H^s|E$ is zero for $\mathcal H^s$-almost every $x\notin E$.
\end{theorem}
\begin{proof}[Formalized in FormalizingGMT]
\leanok
\uses{gmt:hausdorff-regular}
Import the existing declaration.
\end{proof}

\begin{definition}[GMT: vague convergence]
\label{gmt:vague}
\lean{MeasureTheory.Measure.WeaklyConverges}
\leanok
Convergence of a sequence of measures is tested against every real continuous compactly supported function. This does not require finite total mass.
\end{definition}

\begin{definition}[GMT: blow-up map]
\label{gmt:blowup}
\lean{blowUpMap}
\leanok
The blow-up map on Euclidean space is $T_{a,r}(z)=r^{-1}(z-a)$. Geometric use assumes $r>0$.
\end{definition}

\begin{lemma}[GMT: measurable blow-ups and ball identities]
\label{gmt:blowup-balls}
\lean{measurable_blowUpMap,map_blowUp_apply_ball,map_blowUp_apply_closedBall}
\leanok
\uses{gmt:blowup}
Blow-ups are measurable. For $r>0$, the pushforward mass of $B(0,\rho)$ or $\overline B(0,\rho)$ equals the original mass at centre $a$ and radius $r\rho$.
\end{lemma}
\begin{proof}[Formalized in FormalizingGMT]
\leanok
Import the existing declaration.
\end{proof}

\begin{lemma}[GMT: regularity of scaled blow-ups]
\label{gmt:blowup-regular}
\lean{regular_map_blowUp,regular_smul_map_blowUp}
\leanok
\uses{gmt:blowup}
For $r\ne0$, blow-up pushforwards preserve regularity; multiplying by a finite extended nonnegative scalar preserves it too.
\end{lemma}
\begin{proof}[Formalized in FormalizingGMT]
\leanok
Import the existing declaration.
\end{proof}

\begin{definition}[GMT: tangent measures]
\label{def:tangent-api}
\lean{IsTangentMeasure}
\leanok
\uses{gmt:vague,gmt:blowup}
A tangent is a nonzero regular measure obtained as a vague limit of $c_i(T_{a,r_i})_\#\mu$ with $r_i>0$, $r_i\to0$, and $0<c_i<\infty$. The scalar sequence is arbitrary.
\end{definition}

\begin{theorem}[GMT: compact/open portmanteau]
\label{gmt:portmanteau}
\lean{MeasureTheory.Measure.weaklyConverges_iff_compactOpenBounds}
\leanok
\uses{gmt:vague}
For regular Euclidean measures, vague convergence is equivalent to the upper bounds on compact sets and lower bounds on open sets.
\end{theorem}
\begin{proof}[Formalized in FormalizingGMT]
\leanok
Import the existing declaration.
\end{proof}

\begin{theorem}[GMT: convergence on bounded continuity sets]
\label{gmt:continuity-sets}
\lean{MeasureTheory.Measure.weaklyConverges_iff_boundedContinuitySets}
\leanok
\uses{gmt:vague}
For regular Euclidean measures, vague convergence is equivalent to convergence of masses on bounded measurable sets whose boundary has zero limit mass.
\end{theorem}
\begin{proof}[Formalized in FormalizingGMT]
\leanok
\uses{gmt:portmanteau}
Import the existing declaration.
\end{proof}

\begin{theorem}[GMT: compact mass bounds give a subsequence]
\label{thm:compactness}
\lean{MeasureTheory.exists_vaguelyConvergent_subseq_of_compact_bounded}
\leanok
\uses{gmt:vague}
Uniform mass bounds on every compact set give a strictly indexed subsequence converging vaguely to a regular measure. Nonzeroness of the limit must be proved separately.
\end{theorem}
\begin{proof}[Formalized in FormalizingGMT]
\leanok
Import the existing declaration.
\end{proof}

\begin{lemma}[GMT: density bounds imply support and doubling]
\label{gmt:doubling-bounds}
\lean{mem_support_of_ball_lower_bound,limsup_ball_ratio_lt_top_of_ball_bounds_at}
\leanok
\uses{def:density-api}
At a point with positive lower and finite upper power bounds on sufficiently small closed balls, the point belongs to the support and its upper small-scale open-ball doubling ratio is finite.
\end{lemma}
\begin{proof}[Formalized in FormalizingGMT]
\leanok
Import the existing declaration.
\end{proof}

\begin{theorem}[GMT: origin lies in tangent support]
\label{gmt:origin-support}
\lean{zero_mem_support_of_isTangentMeasure}
\leanok
\uses{def:tangent-api}
For a regular measure at a support point with finite upper doubling ratio, every tangent has zero in its support.
\end{theorem}
\begin{proof}[Formalized in FormalizingGMT]
\leanok
\uses{gmt:portmanteau,gmt:blowup-balls}
Import the existing declaration.
\end{proof}

\begin{theorem}[GMT: mass-normalized tangent sequence]
\label{gmt:normalized-tangent}
\lean{exists_normalized_blowUp_weaklyConverges_of_isTangentMeasure}
\leanok
\uses{def:tangent-api}
Under the same support and doubling assumptions, each tangent is the vague limit of a fixed positive finite multiple of the ball-mass-normalized blow-ups along strictly decreasing radii.
\end{theorem}
\begin{proof}[Formalized in FormalizingGMT]
\leanok
\uses{gmt:origin-support,gmt:continuity-sets}
Import the existing declaration.
\end{proof}

\begin{theorem}[GMT: tangent existence under doubling]
\label{thm:doubling}
\lean{exists_subseq_blowUp_weaklyConverges_tangentMeasure}
\leanok
\uses{def:tangent-api}
At a support point of a regular measure with finite upper doubling ratio, every positive sequence of radii tending to zero has a subsequence of ball-mass-normalized blow-ups converging to a nonzero tangent (Mattila 14.3).
\end{theorem}
\begin{proof}[Formalized in FormalizingGMT]
\leanok
\uses{thm:compactness,gmt:blowup-balls,gmt:blowup-regular,gmt:portmanteau}
Import the existing declaration.
\end{proof}

\begin{theorem}[GMT: restriction preserves tangents]
\label{gmt:restriction-tangents}
\lean{isTangentMeasure_restrict_iff}
\leanok
\uses{def:tangent-api}
For Borel $B$, regular $\mu$ and $\mu|B$, and a support point $a$ with $\mu(B(a,r)\setminus B)/\mu(B(a,r))\to0$, restriction to $B$ preserves the tangent set at $a$ (Mattila 14.5).
\end{theorem}
\begin{proof}[Formalized in FormalizingGMT]
\leanok
\uses{gmt:blowup-regular}
Import the existing declaration.
\end{proof}

\begin{theorem}[GMT: density change preserves tangents]
\label{thm:matching}
\lean{isTangentMeasure_iff_ae_of_density}
\leanok
\uses{def:tangent-api}
If regular measures satisfy $\lambda(A)=\int_A\varphi\,d\mu$ on Borel sets for a nonnegative real locally integrable $\varphi$, their tangent sets coincide at $\lambda$-almost every point (Mattila 14.6).
\end{theorem}
\begin{proof}[Formalized in FormalizingGMT]
\leanok
\uses{gmt:blowup-regular}
Import the existing declaration.
\end{proof}

\begin{lemma}[GMT: closed good sets for density bounds]
\label{gmt:goodsets}
\lean{positiveFiniteDensitySet,RatioOfDensities,goodSet,isClosed_goodSet,exists_goodSet_mem}
\leanok
\uses{def:density-api}
Positive finite density-bound points admit a countable stratification by closed good sets with uniform small-radius power bounds. The rational bounds may approximate the lower/upper density ratio.
\end{lemma}
\begin{proof}[Formalized in FormalizingGMT]
\leanok
\uses{gmt:density-bounds}
Import the existing declaration.
\end{proof}

\begin{lemma}[GMT: one null set of bad density points]
\label{gmt:density-points}
\lean{exists_null_of_not_density_point}
\leanok
For a Borel set $B$ and regular Euclidean $\mu$, one null exceptional set can be removed so that every remaining $a\in B$ satisfies, for each $\gamma>0$, $\mu(\overline B(a,r)\setminus B)\leq\gamma\mu(\overline B(a,r))$ eventually as $r\downarrow0$.
\end{lemma}
\begin{proof}[Formalized in FormalizingGMT]
\leanok
Import the existing declaration.
\end{proof}

\begin{theorem}[GMT: transfer bounds at a density point]
\label{gmt:transfer-ball-bounds}
\lean{tangent_ball_bounds_of_density_point}
\leanok
\uses{def:tangent-api}
Uniform small-ball bounds $td\,r^s\leq\mu(\overline B(y,r))\leq d\,r^s$ on a set transfer at its density points to bounds $tc\,r^s\leq\nu(\overline B(z,r))\leq c\,r^s$ at all support centres and all positive radii of every tangent; $c$ is positive and finite.
\end{theorem}
\begin{proof}[Formalized in FormalizingGMT]
\leanok
\uses{gmt:portmanteau,gmt:blowup-balls,gmt:blowup-regular}
Import the existing declaration.
\end{proof}

\begin{theorem}[GMT: tangent bounds from density ratio]
\label{gmt:mattila-14-7-1}
\lean{mattila_14_7_1}
\leanok
\uses{def:tangent-api,def:density-api}
Outside one $\mu$-null set, every tangent at a positive finite density-bound point has support-centred power bounds with ratio $\Theta_*^s/\Theta^{*s}$ (Mattila 14.7(1)).
\end{theorem}
\begin{proof}[Formalized in FormalizingGMT]
\leanok
\uses{gmt:goodsets,gmt:density-points,gmt:transfer-ball-bounds}
Import the existing declaration.
\end{proof}

\begin{theorem}[GMT: uniform support bounds and all tangents]
\label{gmt:mattila-14-7-4}
\lean{mattila_14_7_4}
\leanok
\uses{def:tangent-api}
If one pair of positive power bounds holds at every support point below one fixed radius, then at every support point tangents exist and all of them satisfy the corresponding global support-centred power bounds (Mattila 14.7(4)).
\end{theorem}
\begin{proof}[Formalized in FormalizingGMT]
\leanok
\uses{thm:doubling,gmt:portmanteau,gmt:blowup-balls}
Import the existing declaration.
\end{proof}

\begin{definition}[GMT: uniform measures]
\label{def:uniform-api}
\lean{IsSUniform}
\leanok
A nonzero regular measure is $s$-uniform when $s>0$ and $\nu(\overline B(y,r))=cr^s$ for one $c>0$, all $y\in\spt\nu$ and all $r>0$.
\end{definition}

\begin{definition}[GMT: positive finite existing density set]
\label{gmt:density-exists-set}
\lean{positiveFiniteDensityExistsSet}
\leanok
\uses{def:density-api}
This set is defined using \texttt{HasDensity} together with a positive finite value of the density. It requires an adapter from our liminf/limsup formulation.
\end{definition}

\begin{lemma}[GMT: existing density gives ratio one and bounds]
\label{gmt:uniform-density-adapter}
\lean{SUniformAux.mem_positiveFiniteDensitySet,SUniformAux.ball_bounds_of_mem}
\leanok
\uses{gmt:density-exists-set,gmt:goodsets}
Membership in the positive-finite-existing-density set implies membership in the two-sided density-bound set with ratio one. Such points have positive lower and finite upper small-ball power bounds.
\end{lemma}
\begin{proof}[Formalized in FormalizingGMT]
\leanok
\uses{gmt:density-bounds}
Import the existing declaration.
\end{proof}

\begin{theorem}[GMT: all tangents are uniform]
\label{thm:uniform-upstream}
\lean{mattila_14_9}
\leanok
\uses{def:uniform-api,def:tangent-api,gmt:density-exists-set}
For $s>0$ and regular $\mu$, outside one null set every tangent at every positive finite existing-density point is $s$-uniform and has the origin in its support (Mattila 14.9).
\end{theorem}
\begin{proof}[Formalized in FormalizingGMT]
\leanok
\uses{gmt:mattila-14-7-1,gmt:uniform-density-adapter,gmt:doubling-bounds,gmt:origin-support}
Import the existing declaration.
\end{proof}

\begin{lemma}[GMT: tangents of uniform measures]
\label{gmt:uniform-tangents}
\lean{SUniformAux.exists_isTangentMeasure,SUniformAux.isSUniform_of_isTangentMeasure}
\leanok
\uses{def:uniform-api,def:tangent-api}
At every support point of an $s$-uniform measure tangents exist, and every such tangent is $s$-uniform with zero in its support.
\end{lemma}
\begin{proof}[Formalized in FormalizingGMT]
\leanok
\uses{gmt:mattila-14-7-4,gmt:origin-support}
Import the existing declaration.
\end{proof}

\begin{lemma}[GMT: uniform measure in a half-space]
\label{gmt:uniform-halfspace}
\lean{SUniformAux.exists_halfSpace}
\leanok
\uses{def:uniform-api}
A uniform measure whose support is not the whole ambient space produces a uniform measure supported in a half-space and containing zero in its support. The proof constructs a tangent, but the exported conclusion does not retain that tangent witness.
\end{lemma}
\begin{proof}[Formalized in FormalizingGMT]
\leanok
\uses{gmt:uniform-tangents,gmt:portmanteau}
Import the existing declaration.
\end{proof}

\begin{lemma}[GMT: bounded radial integrals agree]
\label{gmt:uniform-radial}
\lean{SUniformAux.measure_closedBall_eq_of_mem_support,SUniformAux.setIntegral_comp_dist_eq}
\leanok
\uses{def:uniform-api}
For an $s$-uniform measure, same-radius balls centred on the support have equal mass. The integrals of any continuous radial function over those closed balls agree. Passing to global nonnegative Borel radial tests and identifying the Euclidean radial distribution still need adapters.
\end{lemma}
\begin{proof}[Formalized in FormalizingGMT]
\leanok
Import the existing declaration.
\end{proof}

\begin{lemma}[GMT: half-space to hyperplane]
\label{gmt:uniform-hyperplane}
\lean{SUniformAux.exists_hyperplane}
\leanok
\uses{def:uniform-api}
An $s$-uniform measure containing zero in its support and supported in a half-space produces an $s$-uniform measure supported in a hyperplane. The result does not export the lineage of the blow-ups used in its proof.
\end{lemma}
\begin{proof}[Formalized in FormalizingGMT]
\leanok
\uses{gmt:uniform-tangents,gmt:uniform-radial}
Import the existing declaration.
\end{proof}

\begin{lemma}[GMT: transport from a hyperplane]
\label{gmt:hyperplane-transport}
\lean{SUniformAux.exists_isSUniform_of_hyperplane}
\leanok
\uses{def:uniform-api}
An $s$-uniform measure on $\R^{n+1}$ supported in a linear hyperplane yields an $s$-uniform measure on $\R^n$. The proof uses a linear isometry of the hyperplane; the existential interface does not expose it.
\end{lemma}
\begin{proof}[Formalized in FormalizingGMT]
\leanok
Import the existing declaration.
\end{proof}

\begin{lemma}[GMT: lower-dimensional uniform measure]
\label{gmt:uniform-lower-dim}
\lean{SUniformAux.exists_lower_dim}
\leanok
\uses{def:uniform-api}
A uniform measure on $\R^{n+1}$ whose support is not the whole space yields a uniform measure with the same exponent on $\R^n$. This is not yet the moving-centre pseudotangent theorem needed for Part A.
\end{lemma}
\begin{proof}[Formalized in FormalizingGMT]
\leanok
\uses{gmt:uniform-halfspace,gmt:uniform-hyperplane,gmt:hyperplane-transport}
Import the existing declaration.
\end{proof}

\begin{theorem}[GMT: full-dimensional uniform measures are flat]
\label{gmt:uniform-full-dimension}
\lean{mattila_14_11}
\leanok
\uses{def:uniform-api}
An $n$-uniform measure on $\R^n$ is a positive finite scalar multiple of Lebesgue measure (Mattila 14.11). The source definition of uniformity implies $n>0$.
\end{theorem}
\begin{proof}[Formalized in FormalizingGMT]
\leanok
\uses{gmt:uniform-lower-dim}
Import the existing declaration.
\end{proof}

\section{Additional reuse exposed by the proof audit}

\begin{lemma}[GMT: lower growth excludes full support]
\label{gmt:support-not-full}
\lean{support_ne_univ_of_lower_growth}
\leanok
A regular measure on $\R^n$ with a common positive support-centred lower $s$-power bound has non-full support if $s<n$.
\end{lemma}
\begin{proof}[Existing formalization]
\leanok
\uses{def:uniform-api}
Import the linked declaration at the pinned source revision.
\end{proof}

\begin{lemma}[GMT: quadratic barycentre estimate]
\label{gmt:barycenter-quadratic-bound}
\lean{SUniformAux.abs_inner_integral_le}
\leanok
For a uniform measure with zero in its support and each positive radius $r$, the inner product of a support point $y$ with the vector integral over $\overline B(0,r)$ is bounded in absolute value by $C|y|^2$ for $|y|\leq r$.
\end{lemma}
\begin{proof}[Existing formalization]
\leanok
\uses{def:uniform-api,gmt:uniform-radial}
Import the linked declaration at the pinned source revision.
\end{proof}

\begin{lemma}[GMT: tangent support from a quadratic estimate]
\label{gmt:quadratic-support-tangent}
\lean{SUniformAux.support_tangent_subset_of_estimate}
\leanok
If $|\langle y,w\rangle|\leq C|y|^2$ holds on the support near zero, every tangent at zero is supported in $w^\perp$. The estimate has nonnegative finite $C$ and a positive radius of validity.
\end{lemma}
\begin{proof}[Existing formalization]
\leanok
\uses{def:uniform-api,def:tangent-api}
Import the linked declaration at the pinned source revision.
\end{proof}

\begin{lemma}[GMT: vanishing half-space integrals force a hyperplane]
\label{gmt:zero-barycenter-hyperplane}
\lean{SUniformAux.support_subset_of_integral_inner_eq_zero}
\leanok
For a uniform measure supported where $\langle z,e\rangle\geq0$, vanishing of the integral of this inner product over every positive-radius centred closed ball forces its support into $e^\perp$.
\end{lemma}
\begin{proof}[Existing formalization]
\leanok
\uses{def:uniform-api}
Import the linked declaration at the pinned source revision.
\end{proof}

\begin{lemma}[GMT: compact-test estimates and perturbation]
\label{gmt:compact-test-perturbation}
\lean{norm_integral_cc_le,weaklyConverges_of_tendsto_integral_sub}
\leanok
A compact-test integral is bounded by the test supremum norm times the finite mass of a containing support set. Two sequences whose compact-test integral differences tend to zero have the same vague limits.
\end{lemma}
\begin{proof}[Existing formalization]
\leanok
\uses{gmt:vague}
Import the linked declaration at the pinned source revision.
\end{proof}

\begin{lemma}[GMT: negligible discarded mass in normalized blow-ups]
\label{gmt:scaled-discarded-mass}
\lean{tendsto_nhdsGT_zero_of_pos,eventually_smul_measure_closedBall_le,tendsto_mul_measure_diff_zero}
\leanok
Positive radii tending to zero tend through the positive neighborhood filter. Convergent scaled blow-ups have eventual local ball bounds. At a support point where a restriction has relative density one, those bounds make the scaled discarded ball mass tend to zero.
\end{lemma}
\begin{proof}[Existing formalization]
\leanok
\uses{gmt:blowup-balls,gmt:portmanteau}
Import the linked declaration at the pinned source revision.
\end{proof}

\begin{lemma}[mathlib: measurable hull intersection identity]
\label{mathlib:measurable-hull}
\lean{MeasureTheory.Measure.measure_toMeasurable_inter}
\leanok
For a set $D$ of finite measure and measurable $U$, $\mu(\operatorname{toMeasurable}_\mu D\cap U)=\mu(D\cap U)$. No measurability of $D$ is required.
\end{lemma}
\begin{proof}[Existing formalization]
\leanok
Import the linked declaration at the pinned source revision.
\end{proof}

\begin{lemma}[mathlib: disjoint almost-everywhere ball covering]
\label{mathlib:besicovitch-cover}
\lean{Besicovitch.exists_disjoint_closedBall_covering_ae}
\leanok
For an s-finite measure in a space with the Besicovitch covering property, a fine family of positive-radius closed balls admits a countable disjoint subfamily covering the target set almost everywhere, with prescribed positive upper radius bounds.
\end{lemma}
\begin{proof}[Existing formalization]
\leanok
Import the linked declaration at the pinned source revision.
\end{proof}

\begin{lemma}[mathlib: differentiation of measures and restrictions]
\label{mathlib:besicovitch-differentiation}
\lean{Besicovitch.ae_tendsto_rnDeriv,Besicovitch.ae_tendsto_measure_inter_div_of_measurableSet,Besicovitch.ae_tendsto_measure_inter_div}
\leanok
For locally finite measures on finite-dimensional real normed spaces, closed-ball measure ratios converge almost everywhere to the Radon--Nikodym density. Restrictions have relative density one almost everywhere on the restricted set; the arbitrary-set version uses the restricted measure.
\end{lemma}
\begin{proof}[Existing formalization]
\leanok
Import the linked declaration at the pinned source revision.
\end{proof}

\begin{lemma}[mathlib: disjoint selection with enlarged ball covering]
\label{mathlib:vitali-enlargement}
\lean{Vitali.exists_disjoint_subfamily_covering_enlargement_ball}
\leanok
A bounded-radius ball family admits a disjoint subfamily such that each original ball is contained in an enlargement of a selected ball by any fixed factor greater than three. This is distinct from a disjoint almost-everywhere covering without enlargement.
\end{lemma}
\begin{proof}[Existing formalization]
\leanok
Import the linked declaration at the pinned source revision.
\end{proof}



\section{Further generic mathlib infrastructure}
\begin{lemma}[Derivative-controlled closed covers and measurable partitions]
\label{mathlib:derivative-partitions}
\lean{exists_closed_cover_approximatesLinearOn_of_hasFDerivWithinAt, exists_partition_approximatesLinearOn_of_hasFDerivWithinAt}
\leanok
For a differentiable map between finite-dimensional normed spaces, mathlib supplies
countable pieces with uniform linear approximation by actual derivatives.
These generic results apply in codimension; the same-dimensional change-of-variables
formula in that file does not itself supply the graph area formula.
\end{lemma}
\begin{proof}[Existing upstream proof]
\leanok
Source: \texttt{Mathlib/MeasureTheory/Function/Jacobian.lean}.
\end{proof}

\begin{definition}[Intrinsic standard Gaussian measure]
\label{mathlib:standard-gaussian}
\lean{ProbabilityTheory.stdGaussian, ProbabilityTheory.stdGaussian_map, ProbabilityTheory.stdGaussian_eq_map_pi_orthonormalBasis}
\leanok
Mathlib defines the standard Gaussian measure on a finite-dimensional real inner
product space, proves its behavior under linear isometries, and identifies it with
the orthonormal-coordinate product Gaussian. Use this measure on a plane subtype
instead of defining a second Gaussian distribution.
\end{definition}

\begin{lemma}[Uniform measures cannot exceed ambient dimension]
\label{gmt:uniform-dimension-bound}
\lean{SUniformAux.not_isSUniform_of_lt}
\leanok
\uses{def:uniform-api}
There is no $s$-uniform measure on $\mathbb R^n$ when $n<s$.
\end{lemma}
\begin{proof}[Existing upstream proof]
\leanok
FormalizingGMT supplies the obstruction directly. Together with tangent existence
and almost-everywhere uniformity, it proves the excessive-dimension endpoint.
\end{proof}

\begin{lemma}[GMT: closed sets controlling all balls through a point]
\label{gmt:goodballs}
\lean{goodBallSet, isClosed_goodBallSet}
\leanok
\uses{def:density-api}
The set of points for which every sufficiently small closed ball containing
the point satisfies a fixed upper power bound is closed.
\end{lemma}
\begin{proof}[Existing upstream proof]
\leanok
Enlarge a ball slightly, approximate the centre point by good points,
and let the enlargement tend to zero.
\end{proof}

\begin{lemma}[mathlib: mass-distribution bound for Hausdorff measure]
\label{mathlib:hausdorff-mass-bound}
\lean{MeasureTheory.Measure.le_hausdorffMeasure}
\leanok
A measure bounded by the diameter to a real power on every sufficiently
small set is at most the corresponding Hausdorff measure.
\end{lemma}
\begin{proof}[Existing mathlib proof]
\leanok
Use the maximal outer measure defining the metric Hausdorff construction.
\end{proof}

\begin{lemma}[mathlib: compact tests determine regular measures]
\label{mathlib:regular-measure-ext}
\lean{MeasureTheory.Measure.ext_of_integral_eq_on_compactlySupported}
\leanok
On a locally compact Hausdorff space with its Borel measurable structure,
two regular measures are equal if their integrals agree against every
real-valued compactly supported continuous function.
\end{lemma}
\begin{proof}[Existing mathlib proof]
\leanok
Use regularity and continuous compactly supported approximations to compact sets.
\end{proof}
